\chapter{Principio de Ambivalencia: orientaciones conjugadas \texorpdfstring{\(\pi/\varphi^2\)}{π/φ²} y \texorpdfstring{\(\varphi/\pi^2\)}{φ/π²}, involución de hojas y realización euclídea del principio de equivalencia.}
\label{chap:83-16}

\section{Secciones compatibles y teoría de la medición}
\label{p12-83-c16-s1:chap:medida}

La noción dimensional de medida se formula sobre sistemas inversos de estados
compatibles. Antes de asignar una cifra a un estado se especifican su historia
de refinamiento, el acoplamiento con el aparato y la aplicación que produce el
observable. Esta secuencia permite separar la existencia de un valor real de
la operación física mediante la cual ese valor se registra.

\subsection{El número como sección compatible}

Sea \((\mathcal S_n,r_{n+1,n})_{n\geq0}\) un sistema inverso de estados
supervivientes. Una sección compatible es una sucesión
\(x=(x_n)_{n\geq0}\) tal que
\[
 r_{n+1,n}(x_{n+1})=x_n
 \qquad(n\geq0).
\]
Su espacio es el límite inverso
\[
 \mathcal X_\infty=\varprojlim_n\mathcal S_n.
\]
Cuando cada \(x_n\) determina un intervalo \(I_n(x_n)\subset\mathbb R\), las
clausuras \(\overline{I_n(x_n)}\) son compactas no vacías, están anidadas y
sus diámetros tienden a cero, el teorema de los intervalos encajados garantiza
que su intersección contiene un único punto. La evaluación arquimediana se
define entonces por
\[
 \mathcal P_{\rm arq}(x)
 =\text{el único elemento de }
   \bigcap_{n\geq0}\overline{I_n(x_n)}.
\]
El número HMT es la sección compatible
\(x\in\mathcal X_\infty\). El lector \(L\) es estructura adicional, y el
acoplamiento \((x,L)\) constituye su presentación leída o medición. Dos
secciones pueden poseer la misma evaluación real y diferir en sus coordenadas
de cociente, orientación o frontera; la proyección arquimediana no identifica
por definición esas genealogías.

\subsection{Compatibilidad entre cilindros ternarios y milésimos}

Una palabra ternaria de longitud \(L\), cuyo valor entero es \(N\), determina
el cilindro
\[
 I_{N,L}=\left[\frac{N}{3^L},\frac{N+1}{3^L}\right).
\]
Una palabra de \(m\) bloques en base \(1000\), cuyo valor entero es \(D\),
determina
\[
 J_{D,m}=\left[\frac{D}{1000^m},\frac{D+1}{1000^m}\right).
\]

\begin{lemma}[Compatibilidad exacta de cilindros]
Los cilindros \(I_{N,L}\) y \(J_{D,m}\) tienen intersección no vacía si y
sólo si
\[
 N\,1000^m<(D+1)3^L,
 \qquad
 D\,3^L<(N+1)1000^m.
\]
\end{lemma}
\begin{proof}
Dos intervalos semiabiertos \([a,b)\) y \([c,d)\) se intersecan exactamente
cuando \(a<d\) y \(c<b\). Sustituyendo los cuatro extremos y multiplicando
por el entero positivo \(3^L1000^m\) se obtienen las desigualdades. La
decisión se reduce, por tanto, a aritmética entera.
\end{proof}

La prolongación de una palabra refina su cilindro. Una cadena compatible cuyas
clausuras son compactas, no vacías y anidadas, con diámetro convergente a cero,
determina una única evaluación real. El empleo de las clausuras es esencial:
los intervalos semiabiertos anidados pueden tener intersección vacía cuando el
límite coincide con un extremo excluido. El estado completo conserva las
palabras, las transformaciones y los acarreos que producen la evaluación.

\subsection{Acoplamiento reversible entre sistema y aparato}

Sean \(X\) el espacio de estados del sistema, \(M\) el espacio de estados del
aparato y \(m_0\in M\) su preparación inicial. Un acoplamiento es una
aplicación
\[
 U:X\times M\longrightarrow X\times M.
\]
Cuando \(U\) es biyectiva, el estado conjunto conserva la información. Una
aplicación de evaluación \(q_M:M\to Y\) puede ser no inyectiva y publicar
únicamente un observable \(y\in Y\).

\begin{definition}[Medición]
Una medición HMT es la terna
\[
 \mathfrak M=(U,q_M,\mathcal C),
\]
donde \(U\) acopla sistema y aparato, \(q_M\) evalúa el estado del aparato y
\(\mathcal C\subseteq X\times M\) es la condición de compatibilidad. Su
salida sobre \(x\) es
\[
 \operatorname{Med}_{\mathfrak M}(x)
 =q_M\!\left(\operatorname{pr}_M U(x,m_0)\right),
\]
siempre que el par acoplado satisfaga \(\mathcal C\).
\end{definition}

La definición establece una composición precisa de tres operaciones:
acoplamiento, restricción al dominio compatible y evaluación del observable.

\begin{theorem}[Localización de la pérdida de información]
Sea
\[
 \Omega_M:=q_M\circ\operatorname{pr}_M:X\times M\longrightarrow Y.
\]
Si \(U\) es biyectiva y el estado final completo se conserva, las clases de
entradas indistinguibles por el observable publicado son exactamente las
preimágenes por \(U\) de las fibras de \(\Omega_M\), salvo que se componga
además con una eliminación o reinicialización posterior del estado.
\end{theorem}
\begin{proof}
La biyectividad de \(U\) permite recuperar \((x,m_0)\) desde el estado final
conjunto. Para dos entradas \(z,z'\) del dominio compatible, la igualdad de
registros equivale a
\(\Omega_M(Uz)=\Omega_M(Uz')\). Por tanto, \(Uz\) y \(Uz'\) pertenecen a una
misma fibra de la aplicación compuesta \(q_M\circ\operatorname{pr}_M\). Esta
composición incluye tanto la pérdida debida a proyectar sobre el aparato como
la debida a evaluar su estado. Cualquier borrado posterior introduce una
aplicación adicional y debe analizarse por separado.
\end{proof}

\begin{figure}[htbp]
\centering
\resizebox{\textwidth}{!}{%
\begin{tikzpicture}[node distance=13mm and 15mm,>=Latex,
box/.style={draw=hmtblue,rounded corners,fill=hmtlight,align=center,minimum height=9mm,minimum width=29mm}]
\node[box] (s) {estado completo\\[-1mm]$x$};
\node[box,right=of s] (u) {acoplamiento\\[-1mm]$U(x,m_0)$};
\node[box,right=of u] (c) {restricción\\[-1mm]$\mathcal C$};
\node[box,right=of c] (q) {evaluación\\[-1mm]$q_M$};
\node[box,right=of q] (y) {observable\\[-1mm]$y$};
\draw[->,thick] (s)--(u); \draw[->,thick] (u)--(c);
\draw[->,thick] (c)--(q); \draw[->,thick] (q)--(y);
\end{tikzpicture}
}
\caption{Composición de una medición. La aplicación de evaluación actúa
después del acoplamiento y de la restricción de compatibilidad; el estado
conjunto permanece disponible para la transformación inversa.}
\label{p12-83-c16-s1:fig:medicion}
\end{figure}

\subsection{Principio de Ambivalencia: involución bidireccional de los lectores π/φ² y φ/π²}

Sean
\[
 q_A:\mathcal X_{\rm TPK}\to\mathcal X_A,
 \qquad
 q_V:\mathcal X_{\rm TPK}\to\mathcal X_V
\]
dos aplicaciones cociente asociadas, respectivamente, a los observables de
apertura y retención, y sea
\(T:\mathcal X_{\rm TPK}\to\mathcal M\) la variable de memoria
transportada.

\begin{principle}[Ambivalencia]
El estado completo precede a sus proyecciones observables. Las relaciones de
equivalencia inducidas por \(q_A\) y \(q_V\) no se suponen idénticas, y la
compatibilidad entre ambas se registra mediante \(T\). El mapa conjunto
\[
 (q_A,q_V,T):\mathcal X_{\rm TPK}
 \longrightarrow\mathcal X_A\times\mathcal X_V\times\mathcal M
\]
se adopta como condición de separación de los estados físicamente
distinguibles, módulo el grupo de calibre declarado.
\end{principle}

\begin{proposition}[Reconstrucción bajo la condición de separación]
Si el mapa conjunto anterior es inyectivo módulo calibre, cada clase de
estado completo queda determinada unívocamente por sus dos observables y su
variable de memoria.
\end{proposition}
\begin{proof}
Dos estados con la misma terna de imágenes pertenecen a una misma clase de
calibre por la hipótesis de inyectividad. Cada fibra del mapa conjunto
contiene, por tanto, a lo sumo una clase.
\end{proof}

El observable \(q_A\) registra las prolongaciones accesibles y \(q_V\) la
componente retenida. Una interpretación material de la retención exige un
observable de masa que factorice por \(q_V\); no se deduce únicamente del
principio de separación.

\subsection{Determinación relacional y orientación temporal}

Sea \([\gamma]_{\mathcal G}\) la clase de una trayectoria en el grupoide de
estados y sea \(\beta_\gamma\) su dato terminal de frontera.
\begin{principle}[Determinación relacional]
Un observable local admisible tiene la forma
\[
 \mathcal O_{\rm loc}(x;\gamma)
 =F\bigl(x,[\gamma]_{\mathcal G},\beta_\gamma\bigr)
\]
y es invariante bajo las equivalencias de trayectoria que preservan el
registro ordenado. En consecuencia, la coordenada puntual visible no determina
por sí sola el observable cuando la evaluación ha identificado clases de
transporte diferentes.
\end{principle}

La formulación es relacional en un sentido preciso: el argumento de
\(\mathcal O_{\rm loc}\) contiene el estado, la clase global de transporte y
la frontera. Éste es el dominio sobre el que puede compararse posteriormente
la teoría con la tradición asociada al principio de Mach. No se presupone que
toda propiedad inercial haya sido derivada; se especifica el tipo del mapa que
debe realizarla.

Si la actualización sobre \(\mathcal X_{\rm TPK}\) es biyectiva y el registro
conserva las coordenadas requeridas por su inversa, la dinámica completa es
reversible. Una orientación temporal observable aparece al componerla con un
lector no inyectivo, reinicializar un dato de frontera o seleccionar una sola
prolongación entre varias. La irreversibilidad se localiza, por tanto, en un
morfismo definido de pérdida o selección de información, no en la mera
ordenación de las etiquetas.

\subsection{Relación con los principios variacionales clásicos}

El principio de Hamilton exige que la primera variación de una acción se anule
en la trayectoria física; no afirma, en general, que esa trayectoria sea un
mínimo global. El funcional discreto de interacción que se construye en el
\cref{chap:pantalla}, en cambio, ordena una familia finita de rutas por el
número de cambios de orientación y de ley aritmética. Ambos principios
coinciden sólo después de construir un límite en el que el coste discreto
induzca una acción y su primera variación reproduzca las ecuaciones de
movimiento. Esta distinción impide sustituir una correspondencia futura por
una identidad de nombres.

\subsection{Condiciones de una evaluación prospectiva}

Una evaluación prospectiva queda matemáticamente determinada cuando, antes de
consultar su salida, se fijan el espacio de estados, el acoplamiento, las
aplicaciones cociente, la variable de memoria, la regla de compatibilidad y el
criterio de contraste. Si sobreviven varias secciones, la salida es la familia
completa. La asignación de frecuencias requiere entonces una medida de
probabilidad o una dinámica adicional explícita.


\subsection{Cociente de modos y correspondencia areal--volumétrica}
\label{p12-83-c16-s2:sec:cociente-modos-fav}

El signo del calendario, la etiqueta de modo y la hoja geométrica son tres
objetos tipados. En particular, no se identifica el signo de fase con el
observable angular sobre dígitos. Para la correspondencia geométrica fijamos el
lector HMT tipado
\[
 \mathfrak g_{\rm av}(\Sigma)=\mathscr H_{\rm ar},
 \qquad
 \mathfrak g_{\rm av}(\Pi)=\mathscr H_{\rm vol}.
\]
La primera asignación expresa la lectura aditiva o de borde; la segunda, la
lectura multiplicativa o de interior. El mapa \(\mathfrak g_{\rm av}\) es
parte del tipado de hojas. El resultado siguiente demuestra que, una vez
fijado ese tipado, la dinámica TPK fuerza la incidencia entre ellas.

\begin{theorem}[Descenso necesario de la incidencia de hojas]
\label{p12-83-c16-s2:thm:descenso-necesario-fav}
Sea \(m_r\) el modo obtenido después de aplicar la fase \(r\) del bloque
primitivo \((+1,-1,0)\), con la regla
\[
 +1:\ m\mapsto\Sigma,\qquad
 -1:\ m\mapsto\Pi,\qquad
 0:\ m\mapsto m.
\]
La órbita cíclica de modos es \((\Sigma,\Pi,\Pi)\). Si
\[
 N_3(i,j)
 :=
 \#\{r\in\mathbb Z/3\mathbb Z:(m_r,m_{r+1})=(i,j)\},
\]
entonces, en el orden \((\Sigma,\Pi)\),
\[
 \boxed{
 N_3=
 \begin{pmatrix}0&1\\1&1\end{pmatrix}.}
\]
Por repetición del calendario,
\[
 N_9=3N_3,\qquad N_{27}=9N_3,\qquad N_{108}=36N_3.
\]
Después de aplicar \(\mathfrak g_{\rm av}\), la matriz primitiva de
incidencia es, en el orden
\((\mathscr H_{\rm ar},\mathscr H_{\rm vol})\),
\[
 \boxed{
 F_{\rm av}=
 \begin{pmatrix}0&1\\1&1\end{pmatrix}.}
\]
Así, las dos transiciones cruzadas, la única autorretención volumétrica, la
ausencia de autorretención areal y las multiplicidades primitivas unitarias
son consecuencias del cociente de modos de TPK; no son cuatro postulados
añadidos.
\end{theorem}

\begin{proof}
Los tres pares consecutivos, con cierre cíclico, son
\[
 \Sigma\longrightarrow\Pi,\qquad
 \Pi\longrightarrow\Pi,\qquad
 \Pi\longrightarrow\Sigma.
\]
Cada uno aparece una vez y el par
\(\Sigma\to\Sigma\) no aparece. Esto determina las cuatro entradas de
\(N_3\). Los calendarios de longitudes \(9,27,108\) contienen,
respectivamente, \(3,9,36\) copias del bloque primitivo, por lo que sus
matrices de incidencia son los múltiplos indicados. El máximo común divisor
de las entradas no nulas de cada matriz larga elimina únicamente esa
repetición global y devuelve \(N_3\). Finalmente,
\(\mathfrak g_{\rm av}\) transporta \(\Sigma,\Pi\) a las hojas areal y
volumétrica sin alterar la incidencia.
\end{proof}

La realización hiperbólica real que utiliza esta matriz y prueba su
invariancia lorentziana se formula en
\cref{sec:tres-generadores-concretos-trit}; allí \(F_{\rm av}\) se usa
como antecedente, sin repetir su derivación.

Este teorema debe distinguirse de una afirmación de Markov más fuerte.
Olvidar la fase no es un cociente de Markov fuerte de
\(\mathcal K_{\rm ph}\): el modo visible no determina cuál de los tres
operadores \(I,E_+,E_\times\) actúa en la iteración elemental siguiente. Lo que desciende
sin hipótesis adicional es el multigrafo de aristas de un bloque TPK\@.

Existe, sin embargo, una completación de memoria uno determinada de forma
única por ese descenso. Sea \(X_{\rm av}\) el menor subshift de tipo finito
sobre
\(\{\mathscr H_{\rm ar},\mathscr H_{\rm vol}\}\) que contiene todas las
palabras obtenidas al olvidar la fase y que sólo conserva pares consecutivos.
Todo subshift de un paso con esa propiedad debe admitir exactamente los tres
pares observados; el menor de ellos excluye el único par no observado. Por
tanto su matriz de adyacencia es \(F_{\rm av}\).

\begin{corollary}[Medida de Parry del envolvente de memoria uno]
\label{p12-83-c16-s2:cor:parry-envolvente-fav}
El subshift \(X_{\rm av}\) es primitivo, tiene entropía topológica
\(\log\varphi\) y posee una única medida de máxima entropía. Sus pesos
estacionarios de vértice satisfacen
\[
 \boxed{
 \frac{\mu_{\rm ar}}{\mu_{\rm vol}}=\varphi^{-2}.}
\]
\end{corollary}

\begin{proof}
La raíz de Perron de \(F_{\rm av}\) es \(\varphi\), y un vector de Perron
positivo es \((1,\varphi)\). Como la matriz es simétrica, los pesos de vértice
de Parry son proporcionales al producto de los vectores izquierdo y derecho,
es decir, a \((1,\varphi^2)\).
\end{proof}

La medida de Parry pertenece al envolvente de caminos de memoria uno. No es
la imagen de \(\tau_{468}\), ni la frecuencia temporal de la órbita periódica
de modos. Esta última es
\[
 \nu_{\rm cyc}(\mathscr H_{\rm ar})=\frac13,\qquad
 \nu_{\rm cyc}(\mathscr H_{\rm vol})=\frac23,
\]
mientras \(\tau_{468}\) es uniforme sobre emisiones y la medida de Parry
pondera historias de renovación de longitud arbitraria. Las tres medidas
responden a preguntas distintas. La razón irracional
\(\varphi^{-2}\) queda así derivada canónicamente del envolvente de caminos,
sin atribuirla a un cociente de cardinales finitos.

\subsubsection{Del conteo cronológico al retorno marcado por \texorpdfstring{\(\pi\)}{π}}
\label{p12-83-c16-s2:sec:retorno-marcado-pi-phi2}

Las dos fórmulas
\[
 \left(\frac13,\frac23\right)
 \qquad\text{y}\qquad
 \frac{\mu_{\rm ar}}{\mu_{\rm vol}}=\varphi^{-2}
\]
no son aproximaciones una de otra. La primera cuenta los tres instantes del
calendario primitivo \((\Sigma,\Pi,\Pi)\); la segunda pondera todas las
historias admisibles del envolvente de memoria uno. Por eso la frecuencia
cronológica es racional y la razón de Parry es irracional. En lo que sigue
escribimos
\[
 \frac{\pi_{\rm HMT}}{\varphi^2}
 =\pi_{\rm HMT}\varphi^{-2};
\]
no se trata de dividir por \(\varphi^{-2}\).

La aparición de \(\varphi^{-2}\) puede verse directamente en la dinámica
lineal. Si \(F_n\) denota el \(n\)-ésimo número de Fibonacci,
\[
 F_{\rm av}^{\,n}
 =
 \begin{pmatrix}
  F_{n-1}&F_n\\
  F_n&F_{n+1}
 \end{pmatrix}
 \qquad(n\geq1).
\]
Los autovalores de \(F_{\rm av}\) son
\(\varphi\) y \(-\varphi^{-1}\). La rama contraída invierte orientación en
una iteración. Al pasar a la memoria cuadrática, es decir, a
\(\operatorname{Sym}^2(\mathbb R^2)\), esa inversión se registra antes de
desaparecer el signo. En la base \((x^2,2xy,y^2)\), representando por
filas las imágenes de los vectores básicos, el operador inducido es
\[
 \operatorname{Sym}^2(F_{\rm av})=
 \begin{pmatrix}
  0&0&1\\
  0&1&2\\
  1&1&1
 \end{pmatrix},
 \qquad
 \chi(t)=(t+1)(t^2-3t+1).
\]
Sus tres autovalores son
\[
 \varphi^2,\qquad -1,\qquad \varphi^{-2}.
\]
Así, \(\varphi^{-2}\) es el modo estable positivo de la memoria de segundo
orden. Este paso explica por qué la razón áurea aparece aquí como ley de
retorno y no como media estadística añadida.

La clausura \(\pi\) pertenece a otro nivel del mismo estado HMT: es la salida
coinductiva del lector regional ya construido, no un autovalor nuevo de
\(F_{\rm av}\). Sean
\[
 P_{\rm ar}=
 \begin{pmatrix}1&0\\0&0\end{pmatrix},
 \qquad
 P_{\rm vol}=
 \begin{pmatrix}0&0\\0&1\end{pmatrix}.
\]
Entre los operadores positivos que preservan las dos hojas, hay uno y sólo
uno cuya normalización volumétrica es uno y cuya marca areal es la clausura
generada \(\pi_{\rm HMT}\):
\[
 D_\pi=\pi_{\rm HMT}P_{\rm ar}+P_{\rm vol}.
\]
La unicidad es relativa a estas tres cláusulas explícitas: positividad,
diagonalidad respecto de las hojas y las dos normalizaciones. Con
\(e_{\rm ar}=(1,0)^T\) y \(e_{\rm vol}=(0,1)^T\), el retorno marcado de
longitud \(n\) es
\[
 \Theta_n
 :=
 \frac{\langle e_{\rm ar},
 D_\pi F_{\rm av}^{\,n}e_{\rm ar}\rangle}
 {\langle e_{\rm vol},
 F_{\rm av}^{\,n}e_{\rm vol}\rangle}
 =
 \pi_{\rm HMT}\frac{F_{n-1}}{F_{n+1}}.
\]
Por tanto,
\[
 \boxed{\displaystyle
 \lim_{n\to\infty}\Theta_n
 =\frac{\pi_{\rm HMT}}{\varphi^2}.}
\]
La marca \(\pi_{\rm HMT}\) se aplica una sola vez, en el retorno de frontera;
no se reinyecta en cada transición. La combinatoria finita produce el modo
estable \(\varphi^{-2}\), el lector coinductivo produce la clausura
\(\pi_{\rm HMT}\), y \(D_\pi\) compone ambos datos sin invertir su orden
causal.

La longitud nativa \(108\) ofrece un testigo exacto de esa convergencia:
\[
 (F_{\rm av}^{108})_{\rm ar,ar}
 =F_{107}=10284720757613717413913,
\]
\[
 (F_{\rm av}^{108})_{\rm vol,vol}
 =F_{109}=26925748508234281076009.
\]
En consecuencia,
\[
 \left|
 \frac{F_{107}}{F_{109}}-\varphi^{-2}
 \right|
 =
 6.1684979916614407936\ldots\times10^{-46},
\]
y, después de la única marca de clausura,
\[
 \left|
 \pi_{\rm HMT}\frac{F_{107}}{F_{109}}
 -\frac{\pi_{\rm HMT}}{\varphi^2}
 \right|
 =
 1.9378907974286976068\ldots\times10^{-45}.
\]
Estos errores miden la aproximación finita del retorno; no son ajustes de
\(\pi\) ni de \(\varphi\).

\paragraph{Lector cilíndrico de la razón.}
Sean \(I_n^\pi=[\underline\pi_n,\overline\pi_n]\) e
\(I_n^\varphi=[\underline\varphi_n,\overline\varphi_n]\) los cilindros
positivos y encajados producidos por los dos lectores HMT\@. La operación de
intervalos
\[
 \mathcal Q(I,J)
 :=
 \left[
 \frac{\inf I}{(\sup J)^2},
 \frac{\sup I}{(\inf J)^2}
 \right]
\]
es el menor intervalo que contiene \(x/y^2\) para todo
\((x,y)\in I\times J\). Preserva el encaje: si los cilindros de entrada se
refinan, también se refina \(\mathcal Q(I_n^\pi,I_n^\varphi)\). Como sus
diámetros tienden a cero y el límite de \(I_n^\varphi\) es positivo,
\[
 \bigcap_n\mathcal Q(I_n^\pi,I_n^\varphi)
 =
 \left\{\frac{\pi_{\rm HMT}}{\varphi_{\rm HMT}^2}\right\}.
\]
Con los doce bloques en base mil ya publicados, el intervalo cociente fuerza
treinta y cinco cifras decimales:
\[
 \frac{\pi_{\rm HMT}}{\varphi_{\rm HMT}^2}
 =
 1.19998161486432666111577775331680692\ldots.
\]
Esta construcción conserva la procedencia de cada cifra: el numerador
procede del cilindro de clausura y el denominador de la autoescala.

\paragraph{Separación algebraica entre lectura areal, lectura volumétrica y sus codominios respectivos}
La razón completa no puede salir de la sola álgebra racional de
\(F_{\rm av}\). En efecto, sus funciones racionales espectrales pertenecen a
\(\mathbb Q(\sqrt5)\); allí se encuentra \(\varphi^{-2}\), pero no el número
trascendente \(\pi\). Por consiguiente, toda derivación que pretendiese
obtener \(\pi/\varphi^2\) únicamente de la matriz \(2\times2\) estaría
mal tipada. El dato trascendente debe proceder del lector coinductivo de
clausura, o de un cociclo de frontera equivalente demostrado
independientemente. Este control negativo fija la división de trabajo entre
la correspondencia finita y el lector de número.

\subsubsection{Involución de intercambio del cociente}
\label{p12-83-c16-s2:sec:espejo-pi-phi-electron-hmt}

La correspondencia admite dos orientaciones exactas de una misma ley. Si
\[
 \mathscr R(x,y)=\frac{x}{y^2},
 \qquad
 \mathsf J(x,y)=(y,x),
\]
entonces \(\mathsf J^2=I\) y
\[
 \mathscr R(\pi,\varphi)=\frac{\pi}{\varphi^2},
 \qquad
 (\mathscr R\circ\mathsf J)(\pi,\varphi)
 =\frac{\varphi}{\pi^2}.
\]
Se verifican además las identidades
\[
 \mathscr R(\pi,\varphi)\mathscr R(\varphi,\pi)
 =\frac1{\pi\varphi},
 \qquad
 \frac{\mathscr R(\pi,\varphi)}
 {\mathscr R(\varphi,\pi)}
 =\left(\frac{\pi}{\varphi}\right)^3.
\]
Escribiendo
\(\mathcal B_{\rm av}^{(\pi)}=\mathscr R(\pi,\varphi)\), se obtiene
también la reparametrización exacta
\[
 \frac{\varphi}{\pi^2}
 =
 \frac1{\varphi^3
 \bigl(\mathcal B_{\rm av}^{(\pi)}\bigr)^2},
 \qquad
 e^{\varphi/\pi^2}
 =
 \exp\!\left[
 \frac1{\varphi^3
 \bigl(\mathcal B_{\rm av}^{(\pi)}\bigr)^2}
 \right].
\]
Así, el cambio de orientación entre cierre y crecimiento transporta la
razón de frontera a su lectura intercambiada:
\[
 \frac{\pi}{\varphi^2}
 \xleftrightarrow{\ \mathsf J\ }
 \frac{\varphi}{\pi^2}.
\]
La primera orientación lee el cierre areal marcado sobre la memoria
volumétrica cuadrática; la segunda lee el crecimiento interior sobre el
cierre cuadrático. Esta correspondencia es un teorema del lector HMT. La
Mecánica Dimensional aplica después la segunda orientación dentro de la
construcción electrónica, sin convertir al electrón en antecedente de
\(\pi\), \(\varphi\) o de la correspondencia.

\subsubsection{Principio de Ambivalencia de los lectores areal y volumétrico}
\label{p12-83-c16-s2:sec:ambivalencia-grav-inercial-pi-phi2}

El Principio de Ambivalencia afirma que la lectura areal y la lectura
volumétrica son dos funcionales coordinados, pero no intercambiables, de un
mismo estado con memoria. Su coincidencia en la diagonal es
\[
 \mathcal L_{\rm ar}^{\rm diag}(x)
 =\mathcal L_{\rm vol}^{\rm diag}(x).
\]
La igualdad diagonal no borra la ambivalencia fuera de ella.

Sobre la frontera areal--volumétrica, ambos lectores factorizan a través de
una misma amplitud positiva de memoria \(\mathcal M(x)\):
\[
 \mathcal L_{\rm vol}(x)=\mu_{\rm vol}\mathcal M(x),
 \qquad
 \mathcal L_{\rm ar}(x)=
 \pi_{\rm HMT}\mu_{\rm ar}\mathcal M(x).
\]
Entonces el teorema de Parry y la marca de retorno dan
\[
 \boxed{\displaystyle
 \frac{\mathcal L_{\rm ar}(x)}
 {\mathcal L_{\rm vol}(x)}
 =\frac{\pi_{\rm HMT}}{\varphi^2}.}
\]
Con la normalización volumétrica, la diferencia relativa de lectores es
\[
 \frac{\mathcal L_{\rm ar}-\mathcal L_{\rm vol}}
 {\mathcal L_{\rm vol}}
 =
 \frac{\pi_{\rm HMT}}{\varphi^2}-1
 =
 0.1999816148643266611\ldots.
\]
La lectura física de esta correspondencia se realiza más tarde, en el capítulo de
Cartan--Holst, después de construir holonomía, co-tétrada, conexión, torsión
y curvatura. Aquí el resultado es enteramente HMT: una razón exacta entre
dos lectores de la misma memoria de hoja.

Conviene distinguirlos del extractor enriquecido que se utiliza después en
Mecánica Dimensional. Allí una historia completa conserva el conteo de acción
\(n_A\), los doce eventos orientados \(e_0,\ldots,e_{11}\) y la cocadena
torsional a nivel de paso; de ellos se obtienen
\[
 s_C=\sum_{i=0}^{5}(e_i+e_{i+6}),
 \qquad
 W_\pm=6n_A\pm s_C.
\]
Los canales ya generados por HMT satisfacen entonces la identidad exacta
\[
 \exp\!\left(n_AA_{\rm rad}+\frac{s_C}{6}C^\ast_{\rm rad}\right)
 =\left(q_+^{-W_+}q_-^{-W_-}\right)^{1/12}.
\]
Por tanto, el paso ruta--carácter angular no está ausente: se realiza sobre
el registro dodecafásico enriquecido. Lo que este cálculo de fase prueba por sí
solo son \(\Omega_{\rm sw},S,G\); no los sustituye ni permite reconstruir el
registro completo después de haberlo proyectado.

La medida uniforme sobre las \(104\,976\) presentaciones desciende a átomos
de masas \(1/729\), \(1/243\) y \(1/54\), pero no a \(1/468\). En la
microfibra quíntuple de \(\pi\), las dos probabilidades son
\[
 \tau_{468}(\mathscr R_\pi)=\frac5{468},
 \qquad
 F_\#\mu_{\rm seed}(\mathscr R_\pi)=\frac7{729}.
\]
Tanto la elevación promediada como la elevación de fase realizan la primera: no
empujan la medida uniforme sobre semillas, sino que asignan la misma masa
total a cada clase observable. La elevación de fase conserva además la
quietud microscópica en los pasos neutros.

Para relacionar la elevación de fase con las historias enriquecidas de
profundidad seis hace falta un núcleo \(\mathcal K_6\) sobre
\(\mathcal H_6^{\rm cat}\times C_9\). Éste desciende a
\(\widehat{\mathcal K}_{\rm ph}\) a través de
\(\rho_6\times\mathrm{id}_{C_9}\) si y sólo si satisface, en cada fase, la
condición de agregabilidad fuerte
\[
 \sum_{\rho_6(g)=y}\mathcal K_6((h,r),(g,r+1))
 =L_{\tau(r)}(\rho_6(h),y),
\]
independientemente de \(h\) dentro de cada fibra de \(\rho_6\). En ese caso,
la composición con \(F\) proyecta a \(\mathcal K_{\rm ph}\). La construcción
de una familia \((\mathcal K_m)_{m\geq6}\) que satisfaga esta identidad y
conmute con todos los truncamientos es el paso restante hacia una
transferencia del límite inverso.


\section{Correspondencia areal–volumétrica y razón $\pi /\varphi ^2$}
\label{p12-83-c16-s3:sec:bisagra-hojas-pi-phi2}

La sección anterior construyó la involución dodecafásica, su elevación a
secciones persistentes, la banda de Möbius relativa a una holonomía negativa
y la carta biangular \((A,C^*)\leftrightarrow(q_+,q_-)\). Falta ahora
expresar con precisión cómo se cuentan los retornos de las hojas areal y
volumétrica y por qué la razón áurea aparece al cuadrado.

\subsection{Clases, soportes y hojas son objetos diferentes}

Las clases nonádicas de movimiento son
\[
 C_1=\{1,4,7\},\qquad C_2=\{2,5,8\},\qquad C_0=\{3,6,9\}.
\]
Los soportes empleados por el observable angular son
\[
 \mathcal D_{\rm act}=\{1,3,5,7\},\qquad
 \mathcal D_{\rm evt}=\{2,4,6,8\},\qquad
 \mathcal D_9=\{9\}.
\]
Su cruce exacto es
\[
 B_{\rm na}=\bigl(|C_i\cap D_j|\bigr)_{i,j}
 =\begin{pmatrix}2&1&0\\1&2&0\\1&1&1\end{pmatrix}.
\]

\begin{theorem}[Incidencia triádica]
\label{p12-83-c16-s3:thm:incidencia-triadica-corta}
Se cumplen las identidades
\[
 \det B_{\rm na}=3,
 \qquad
 \operatorname{SNF}(B_{\rm na})=\operatorname{diag}(1,1,3),
\]
\[
 \operatorname{im}B_{\rm na}
 =\{(u,v,w)\in\mathbb Z^3:u+v\equiv0\pmod3\}.
\]
El intercambio simultáneo
\(C_1\leftrightarrow C_2\),
\(\mathcal D_{\rm act}\leftrightarrow\mathcal D_{\rm evt}\) deja fija la
incidencia.
\end{theorem}

\Needspace{5\baselineskip}
El observable del registro
\[
 j_2(d)=
 \begin{cases}
 -1,&d\in\mathcal D_{\rm act},\\
 +1,&d\in\mathcal D_{\rm evt},\\
 0,&d\in\mathcal D_9
 \end{cases}
\]
se compone con el lector geométrico activo:
\[
 -1\mapsto\mathscr H_{\rm ar},\qquad
 +1\mapsto\mathscr H_{\rm vol},\qquad
 0\mapsto\mathscr H_{\rm tor}.
\]
Por tanto,
\[
 \mathfrak g_\triangle(\mathcal D_{\rm act})=\mathscr H_{\rm ar},\qquad
 \mathfrak g_\triangle(\mathcal D_{\rm evt})=\mathscr H_{\rm vol},\qquad
 \mathfrak g_\triangle(\mathcal D_9)=\mathscr H_{\rm tor}.
\]
La matriz prueba el cruce entre clases y soportes. El lector
\(\mathfrak g_\triangle\) declara el paso de soportes a hojas. Omitir este
segundo mapa sería identificar cardinales de objetos distintos.

\subsection{Apertura espectral de la banda mediante la fase angular conjugada}

En la carta activa,
\[
 q_+=e^{-A_{\rm rad}-C^*_{\rm rad}},\qquad
 q_-=e^{-A_{\rm rad}+C^*_{\rm rad}},
\]
y la inversión exacta es
\[
 A_{\rm rad}=-\frac12\log(q_+q_-),
 \qquad
 C^*_{\rm rad}=\frac12\log\frac{q_-}{q_+}.
\]
Así, \(A\) es el modo común y \(C^*\) el modo diferencial. Para cada
altura de torre,
\[
 q_-^n-q_+^n
 =2e^{-nA_{\rm rad}}\sinh(nC^*_{\rm rad}).
\]
La holonomía negativa construida en el
apartado~\ref{sec:persistencia-ruta-mobius-familias} fija el pegado
de Möbius; la última identidad demuestra que \(C^*\) mide la apertura
orientada de las dos hojas conjugadas. En la representación física, esas dos
orientaciones se leen como partícula y antipartícula.

En el lector de borde, una vuelta de la base se parametriza por \(2\pi\) y
termina en la orientación opuesta. El lazo cuadrado, parametrizado por
\(4\pi\), restaura la orientación. Por ello, relativamente al levantamiento y al
lector geométrico declarados, \(2\pi\) es cierre areal proyectado y
\(4\pi\) cierre volumétrico de la cubierta orientada. Esta construcción no
confunde ese doble retorno con una representación espinorial; la
representación, cuando se introduce, es un funtor posterior.

\subsection{Descenso de TPK, envolvente de caminos y conteo de lazos}

El calendario TPK y el observable angular \(j_2\) son mapas distintos. El
primero actualiza el modo aritmético
\(m\in\{\Sigma,\Pi\}\); el segundo clasifica soportes de dígitos. La
bisagra emplea el lector HMT tipado
\[
 \eta_{\rm av}(\Sigma)=\mathscr H_{\rm ar},
 \qquad
 \eta_{\rm av}(\Pi)=\mathscr H_{\rm vol}.
\]
Así, \(\Sigma\) es el modo aditivo o de borde y \(\Pi\) el modo
multiplicativo o de interior. Esta declaración no identifica
\(\tau\) con \(j_2\).

En el bloque primitivo de fase
\[
 \tau_3=(+1,-1,0),
\]
TPK actualiza el modo mediante
\[
 +1:\ m\mapsto\Sigma,\qquad
 -1:\ m\mapsto\Pi,\qquad
 0:\ m\mapsto m.
\]
La órbita cíclica es \((\Sigma,\Pi,\Pi)\). Definamos su matriz de
incidencia de aristas por
\[
 N_3(i,j)=
 \#\{r\in\mathbb Z/3\mathbb Z:(m_r,m_{r+1})=(i,j)\}.
\]

\begin{theorem}[Descenso necesario de las cuatro reglas]
\label{p12-83-c16-s3:thm:descenso-necesario-fav-corta}
En el orden \((\Sigma,\Pi)\),
\[
 N_3=
 \begin{pmatrix}0&1\\1&1\end{pmatrix}.
\]
Además,
\[
 N_9=3N_3,\qquad
 N_{27}=9N_3,\qquad
 N_{108}=36N_3.
\]
Por tanto, tras eliminar el factor común de repetición y aplicar
\(\eta_{\rm av}\), el cociente de modos de TPK induce necesariamente
\[
 \boxed{
 F_{\rm av}=S+P_{\rm vol}
 =
 \begin{pmatrix}0&1\\1&1\end{pmatrix}
 }
\]
donde
\[
 S=\begin{pmatrix}0&1\\1&0\end{pmatrix},
 \qquad
 P_{\rm vol}=\begin{pmatrix}0&0\\0&1\end{pmatrix}.
\]
Las dos transiciones cruzadas, la única autorretención volumétrica, la
ausencia de autorretención areal y las multiplicidades primitivas unitarias
son consecuencias de TPK, no premisas añadidas.
\end{theorem}

\begin{proof}
Los tres pares consecutivos del ciclo son
\[
 \Sigma\longrightarrow\Pi,\qquad
 \Pi\longrightarrow\Pi,\qquad
 \Pi\longrightarrow\Sigma.
\]
Cada uno aparece una vez y \(\Sigma\to\Sigma\) no aparece. Esto determina
las cuatro entradas de \(N_3\). Los calendarios de longitudes
\(9,27,108\) contienen \(3,9,36\) bloques primitivos, respectivamente;
de ahí los tres múltiplos. La reducción por el máximo común divisor de las
entradas no nulas elimina sólo la repetición del bloque. El lector
\(\eta_{\rm av}\) cambia el tipo de los vértices, no su incidencia.
\end{proof}

El teorema anterior es un descenso exacto de aristas. No convierte el olvido
de la fase en un cociente de Markov fuerte: una misma emisión visible puede
ser actualizada por \(I,E_+\) o \(E_\times\), según la fase escondida. De
este hecho se desprende una distinción necesaria entre tres medidas:
\[
 \tau_{468},\qquad
 \nu_{\rm cyc},\qquad
 \mu_{\rm Parry}.
\]
La primera es uniforme sobre las \(468\) emisiones y estacionaria para la
transferencia de fase. La segunda es la frecuencia de la órbita periódica,
\[
 \nu_{\rm cyc}(\mathscr H_{\rm ar})=\frac13,\qquad
 \nu_{\rm cyc}(\mathscr H_{\rm vol})=\frac23.
\]
Ninguna de ellas puede producir directamente el cociente irracional
\(\varphi^{-2}\).

El objeto que sí lo produce está determinado por una propiedad universal.
Sea \(X_{\rm av}\) el menor subshift de tipo finito de memoria uno que
contiene las palabras obtenidas al olvidar la fase y conserva todos sus
pares consecutivos. Todo sistema de un paso que contenga esas palabras debe
admitir
\[
 {\rm ar}\to{\rm vol},\qquad
 {\rm vol}\to{\rm vol},\qquad
 {\rm vol}\to{\rm ar}.
\]
El menor excluye el único par ausente, \({\rm ar}\to{\rm ar}\). En
consecuencia, \(X_{\rm av}\) tiene matriz de adyacencia \(F_{\rm av}\).
Esta es la completación simbólica mínima de memoria uno del cociente TPK; no
es una dinámica microscópica nueva.

\begin{theorem}[Ley Fibonacci y medida de Parry]
\label{p12-83-c16-s3:thm:ley-fibonacci-retornos-corta}
El operador \(F_{\rm av}\) es primitivo y
\[
 F_{\rm av}^{\,n}
 =\begin{pmatrix}F_{n-1}&F_n\\F_n&F_{n+1}\end{pmatrix}
 \qquad(n\geq1).
\]
El subshift \(X_{\rm av}\) tiene entropía topológica
\(\log\varphi\) y una única medida invariante de entropía máxima. Para
esta medida,
\[
 \frac{(F_{\rm av}^{\,n})_{\rm ar,ar}}
 {(F_{\rm av}^{\,n})_{\rm vol,vol}}
 =\frac{F_{n-1}}{F_{n+1}}
 \longrightarrow\varphi^{-2},
 \qquad
 \boxed{\frac{\mu_{\rm ar}}{\mu_{\rm vol}}=\varphi^{-2}.}
\]
\end{theorem}

\begin{proof}
El cuadrado de \(F_{\rm av}\) es estrictamente positivo. La fórmula de
potencias se obtiene por inducción mediante
\(F_{n+1}=F_n+F_{n-1}\). Las entradas diagonales cuentan lazos cerrados.
La raíz de Perron es \(\varphi\) y un vector de Perron positivo es
\((1,\varphi)\). Como la matriz es simétrica, los pesos de vértice de
Parry son proporcionales al producto de los vectores izquierdo y derecho,
es decir, a \((1,\varphi^2)\).
\end{proof}

El factor \(\varphi^{-1}\) es la razón de amplitudes de Perron de una
travesía; \(\varphi^{-2}\) es la razón de pesos de un retorno completo.
Así quedan ordenados los dos exponentes sin introducir dos constantes
independientes.

\subsection{Origen espectral del exponente cuadrático del retorno}

La aparición del cuadrado admite una segunda demostración, puramente
espectral, que no depende de interpretar una media ponderada. El polinomio
característico de \(F_{\rm av}\) es
\[
 \chi_{F_{\rm av}}(t)=t^2-t-1,
\]
de modo que sus dos autovalores son
\[
 \lambda_+=\varphi,\qquad \lambda_-=-\varphi^{-1}.
\]
El signo de \(\lambda_-\) registra la inversión de orientación de una
travesía de la hoja estable. Al elevar la acción a la memoria cuadrática
\(\operatorname{Sym}^2(\mathbb R^2)\), los autovalores son los productos
por pares
\[
 \lambda_+^2=\varphi^2,\qquad
 \lambda_+\lambda_-=-1,\qquad
 \lambda_-^2=\varphi^{-2}.
\]
En la base
\((e_{\rm ar}^2,e_{\rm ar}\odot e_{\rm vol},e_{\rm vol}^2)\), escribiendo
por columnas las coordenadas de las imágenes de los tres vectores básicos, la
matriz es
\[
 \operatorname{Sym}^2(F_{\rm av})=
 \begin{pmatrix}
 0&0&1\\
 0&1&2\\
 1&1&1
 \end{pmatrix},
\]
y satisface
\[
 \chi_{\operatorname{Sym}^2(F_{\rm av})}(t)
 =t^3-2t^2-2t+1
 =(t+1)(t^2-3t+1).
\]
Por tanto, el modo estable cambia de orientación con
\(-\varphi^{-1}\) en una travesía y se convierte en el peso positivo
\(\varphi^{-2}\) cuando se lee como memoria de área o como retorno
cuadrático. Ésta es la razón estructural del exponente dos.

\subsection{Cierre de borde y cociente \texorpdfstring{$\pi/\varphi^2$}{pi/phi2}}

El lector coinductivo de \(\pi\) precede a esta construcción. Definimos la
banda de cierre por composición, sin usar \(\pi\) para seleccionar el
transportador:
\[
 \mathcal B_{\rm av}^{(\pi)}(X_\infty)
 :=\Pi_{\rm HMT}(X_\infty)
 \frac{\mu_{\rm ar}}{\mu_{\rm vol}}.
\]

\paragraph{Retorno marcado \(\Theta_n\) y convergencia de la razón areal–volumétrica}
Sean los proyectores de hoja
\[
 P_{\rm ar}=\begin{pmatrix}1&0\\0&0\end{pmatrix},
 \qquad
 P_{\rm vol}=\begin{pmatrix}0&0\\0&1\end{pmatrix}.
\]
El marcador de frontera se define por
\[
 D_\pi:=\Pi_{\rm HMT}(X_\infty)P_{\rm ar}+P_{\rm vol}.
\]
Es el único operador positivo diagonal respecto de la descomposición de
hojas que asigna peso unitario a la hoja volumétrica y el cierre
\(\Pi_{\rm HMT}(X_\infty)\) a la hoja areal. Su canonicidad es, por tanto,
relativa a tres cláusulas explícitas: diagonalidad de hoja, normalización
volumétrica y marcado areal por el lector de clausura ya construido.

Si \(e_{\rm ar},e_{\rm vol}\) son los vectores de hoja, el funcional de
retorno marcado de profundidad \(n\) es
\[
 \Theta_n(X_\infty):=
 \frac{\langle e_{\rm ar},
 D_\pi F_{\rm av}^{\,n}e_{\rm ar}\rangle}
 {\langle e_{\rm vol},
 F_{\rm av}^{\,n}e_{\rm vol}\rangle}.
\]
La ley de potencias anterior da, sin aproximación,
\[
 \boxed{
 \Theta_n(X_\infty)
 =\Pi_{\rm HMT}(X_\infty)\frac{F_{n-1}}{F_{n+1}}.}
\]
El cierre \(\pi\) se aplica una sola vez, al marcar el retorno de borde; no
se reintroduce en cada iteración de \(F_{\rm av}\).

\begin{corollary}[Razón holográfica areal--volumétrica]
\label{p12-83-c16-s3:cor:pi-phi2-corta}
Relativamente al cociente de modos TPK, a su envolvente mínimo de memoria
uno y al cierre HMT de \(\pi\),
\[
 \boxed{\mathcal B_{\rm av}^{(\pi)}
 =\lim_{n\to\infty}\Theta_n
 =\frac{\pi_{\rm HMT}}{\varphi^2}.}
\]
\end{corollary}

\paragraph{Periodicidad nativa de \(108\) fases y cierre del calendario unitario}
En la profundidad dodecafásica \(108\), las dos cuentas diagonales son
\[
 (F_{\rm av}^{108})_{\rm ar,ar}
 =F_{107}=10284720757613717413913,
\]
\[
 (F_{\rm av}^{108})_{\rm vol,vol}
 =F_{109}=26925748508234281076009.
\]
Por consiguiente,
\[
 \frac{F_{107}}{F_{109}}
 =0.381966011250105151795413165634361882\ldots,
\]
y el error respecto de \(\varphi^{-2}\) es
\[
 \left|\frac{F_{107}}{F_{109}}-\varphi^{-2}\right|
 =6.1684979916614407936\ldots\times10^{-46}.
\]
Tras el marcado de borde, el error respecto de
\(\pi/\varphi^2\) es
\[
 \left|\pi\frac{F_{107}}{F_{109}}-\frac{\pi}{\varphi^2}\right|
 =1.9378907974286976068\ldots\times10^{-45}.
\]
Así, \(108\) no se usa para ajustar la razón: es una profundidad nativa del
calendario que ya aproxima el retorno límite con más de cuarenta y cuatro
cifras fraccionarias estables.

\paragraph{Lector intervalar coinductivo: encaje de cilindros y unicidad del límite}
Para intervalos positivos
\[
 I=[a,b],\qquad J=[c,d],\qquad 0<c\leq d,
\]
definimos
\[
 \mathcal Q(I,J):=
 \left[\frac{a}{d^2},\frac{b}{c^2}\right].
\]
La función \((x,y)\mapsto x/y^2\) es creciente en \(x\) y decreciente en
\(y>0\). Por ello \(\mathcal Q(I,J)\) es exactamente el menor intervalo
que contiene todos los cocientes \(x/y^2\) con
\((x,y)\in I\times J\); no es una cota elegida por redondeo.

Sean \(I_n^\pi\) e \(I_n^\varphi\) los cilindros positivos anidados
generados por HMT\@. Entonces
\[
 \mathcal Q(I_{n+1}^\pi,I_{n+1}^\varphi)
 \subseteq
 \mathcal Q(I_n^\pi,I_n^\varphi),
\]
y, como
\[
 \bigcap_n I_n^\pi=\{\pi_{\rm HMT}\},
 \qquad
 \bigcap_n I_n^\varphi=\{\varphi_{\rm HMT}\},
\]
se obtiene
\[
 \bigcap_n\mathcal Q(I_n^\pi,I_n^\varphi)
 =\left\{\frac{\pi_{\rm HMT}}{\varphi_{\rm HMT}^2}\right\}.
\]
Cada bloque de base \(1000\) queda forzado en cuanto el intervalo cociente
entra por completo en un único cilindro triádico de la profundidad
correspondiente.

En particular, los cilindros positivos de doce tríadas para \(\pi\) y
\(\varphi\) producen el intervalo cociente
\[
 \left[
 \frac{p^-}{(f^+)^2},
 \frac{p^+}{(f^-)^2}
 \right].
\]
La evaluación certifica treinta y cinco cifras fraccionarias:
\[
 \frac{\pi}{\varphi^2}
 =1.19998161486432666111577775331680692\ldots,
\]
y once bloques completos:
\[
 199|981|614|864|326|661|115|777|753|316|806.
\]

\paragraph{Teorema de separación algebraica.}
La matriz \(F_{\rm av}\), sus potencias y su álgebra espectral racional
viven en una extensión algebraica de \(\mathbb Q\), contenida en
\(\mathbb Q(\sqrt5)\). En consecuencia, ninguna composición que use sólo
operaciones algebraicas racionales sobre \(F_{\rm av}\) puede producir el
factor trascendente \(\pi/\varphi^2\). El operador finito genera
\(\varphi^{-2}\); el lector coinductivo de frontera genera
\(\pi_{\rm HMT}\); el retorno marcado compone ambos sin confundir sus
procedencias. Este control negativo impide presentar \(F_{\rm av}\) como
generador aislado de \(\pi\).

\subsection{Involución de lectores y sección electrónica}
\label{p12-83-c16-s3:sec:espejo-pi-phi-electron}

La relación recién obtenida aparece junto al electrón de una manera más
precisa que una repetición literal. El invariante electrónico activo del
apartado~\ref{sec:persistencia-ruta-mobius-familias} contiene
\[
 \mathfrak e_0=
 \frac{\sqrt3}{4}
 \left(e^{\varphi/\pi^2}-22\alpha_{\rm HMT}^{3}\right)
 (1+15\Delta_4).
\]
Por tanto, el exponente electrónico es \(\varphi/\pi^2\), mientras que la
bisagra de hojas produce \(\pi/\varphi^2\). Los dos escalares no son
iguales ni recíprocos. Forman las dos orientaciones de una misma ley
bivariada. En efecto, definamos
\[
 \mathscr R(x,y):=\frac{x}{y^2},
 \qquad
 \mathsf J(x,y):=(y,x).
\]
Entonces \(\mathsf J^2=I\) y
\[
 \boxed{
 \mathscr R(\pi,\varphi)=\frac{\pi}{\varphi^2},
 \qquad
 (\mathscr R\circ\mathsf J)(\pi,\varphi)
 =\frac{\varphi}{\pi^2}.}
\]
El intercambio de los dos lectores lleva exactamente el retorno marcado
areal--volumétrico a la semilla exponencial del electrón. Además,
\[
 \mathscr R(\pi,\varphi)\mathscr R(\varphi,\pi)
 =\frac1{\pi\varphi},
 \qquad
 \frac{\mathscr R(\pi,\varphi)}
 {\mathscr R(\varphi,\pi)}
 =\left(\frac{\pi}{\varphi}\right)^3.
\]
Si denotamos por
\(\mathcal B_{\rm av}^{(\pi)}=\pi/\varphi^2\) la razón de hojas
marcada, la misma identidad puede escribirse sin introducir una constante
adicional:
\[
 \boxed{
 \frac{\varphi}{\pi^2}
 =
 \frac{1}
 {\varphi^3\bigl(\mathcal B_{\rm av}^{(\pi)}\bigr)^2},}
 \qquad
 e^{\varphi/\pi^2}
 =
 \exp\!\left[
 \frac{1}
 {\varphi^3\bigl(\mathcal B_{\rm av}^{(\pi)}\bigr)^2}
 \right].
\]
Numéricamente,
\[
 \begin{aligned}
 \frac{\varphi}{\pi^2}
 &=0.163941118913672383599752755745666982\ldots,\\
 e^{\varphi/\pi^2}
 &=1.17814494259630678081160095680888910\ldots.
 \end{aligned}
\]

Esta identidad formaliza el contenido bidireccional sin alterar la fuerza de la
ecuación electrónica. La orientación \((\pi,\varphi)\) lee cierre de borde
por memoria cuadrática; la orientación intercambiada
\((\varphi,\pi)\) lee crecimiento interior por cierre cuadrático y entra
exponenciada en \(\mathfrak e_0\). El teorema demuestra que ambas
expresiones pertenecen a la misma órbita de intercambio. Los coeficientes
\(22\) y \(15\) poseen un generador independiente: el
\cref{p12-83-c15-s2:thm:selector-persistencia-electron-rev7} los obtiene, sin consultar la
masa electrónica, como
\[
 -22=-\bigl|\mathcal K_e\setminus\iota(\mathscr R_\pi)\bigr|,
 \qquad
 +15=\bigl|\mathscr R_\pi\times\mathbb F_3\bigr|.
\]
Su estatuto es \textsc{demostrado con estructura de partida explícita}. La
involución actual transporta esa construcción a la orientación electrónica y
no usa el electrón para obtener \(F_{\rm av}\), \(\pi\) o \(\varphi\). La
última forma es un cambio exacto de coordenadas de la semilla electrónica, no
una derivación alternativa de \(\mathfrak e_0\).

\subsection{Identidad pentádica de Rogers--Ramanujan y reconocimiento del invariante de Perron}

La fracción continua de Rogers--Ramanujan proporciona una realización
analítica del mismo invariante. Con
\(q_A=e^{-A_{\rm rad}}\), \(q_C=e^{-C^*_{\rm rad}}\), se obtiene
\[
 R(q_A)R(q_C)-\varphi^{-2}\simeq-1.01864\times10^{-27},
\]
y con los canales de la descomposición partida,
\[
 R(q_+)R(q_-)-\varphi^{-2}\simeq-1.19458\times10^{-21}.
\]
Las diferencias no son cero. La afirmación exacta es el límite
\[
 q_A^{(m)},q_C^{(m)}\to1^-
 \quad\Longrightarrow\quad
 R(q_A^{(m)})R(q_C^{(m)})\to\varphi^{-2}.
\]
Rogers--Ramanujan reconoce así el invariante de Perron en el lector
pentádico; no selecciona por sí sola las coordenadas angulares \(A\), \(C^*\) ni
\(F_{\rm av}\). El antiguo kernel finito \(90/120\) conserva un residuo
aproximado \(7.81\times10^{-9}\) y no se usa como identidad exacta.

\subsection{Involución de los lectores inercial y gravitatorio y reconstrucción bidireccional de la magnitud}

El principio de Ambivalencia distingue el lector de apertura
\(q_A\), el lector de retención \(q_V\) y la memoria transportada. La
bisagra anterior permite formular con exactitud la intuición física, sin
confundirla con el teorema combinatorio. Sea \(\mathcal M(x)\) un carácter
común de memoria y definamos, en la interfaz areal--volumétrica,
\[
 \mathcal I_{\rm av}(x)
 :=\mu_{\rm vol}\mathcal M(x),
 \qquad
 \mathcal G_{\rm av}(x)
 :=\pi_{\rm HMT}\mu_{\rm ar}\mathcal M(x).
\]
Entonces, siempre que \(\mathcal M(x)\ne0\),
\[
 \boxed{
 \frac{\mathcal G_{\rm av}(x)}{\mathcal I_{\rm av}(x)}
 =\frac{\pi_{\rm HMT}}{\varphi^2},}
 \qquad
 \frac{\mathcal G_{\rm av}-\mathcal I_{\rm av}}
 {\mathcal I_{\rm av}}
 =\frac{\pi_{\rm HMT}}{\varphi^2}-1.
\]
La primera fórmula expresa la razón orientada entre los lectores
gravitatorio e inercial; la segunda, su defecto relativo. La L gnomónica
porta la carta euclídea local sobre la que ambos lectores descienden a una
normalización común,
\[
 \frac{\mathcal G_{L}}{\mathcal I_{L}}=1.
\]

La lectura física es estratificada. En la carta euclídea de la L, la
conexión torsional local se proyecta sobre una normalización única y el
principio de equivalencia se realiza como cociente uno. En el estado
enriquecido anterior a esa proyección, la apertura areal y la retención
volumétrica permanecen separadas: el Principio de Ambivalencia conserva sus
dos funcionales y fija el cociente interno \(\pi/\varphi^2\). La equivalencia
local y la ambivalencia de frontera son, así, dos lecturas sucesivas del
mismo antecedente, con dominios y mapas declarados.

La realización física adopta la factorización del observable inercial por la
retención volumétrica y la del observable gravitatorio por la apertura areal
con clausura de borde \(\pi\); ambos comparten el carácter \(\mathcal M\). La
matemática interna determina la correspondencia y su cociente, y el funtor de
Mecánica Dimensional publica sus cartas observables. Una ruptura de cualquiera
de esas factorizaciones constituye el falsador explícito de esta realización.

\paragraph{Frontera cosmológica.}
La hipótesis de que materia oscura o energía oscura sean lecturas efectivas
de cruces no torsionales se conserva únicamente como
una hipótesis de prolongación física. Para convertirla en una proposición
contrastable
deben construirse, al menos, un mapa desde interfaces de hoja a una
geometría cosmológica, los observables que representen densidad y presión y
una comparación prospectiva con datos sin usar esos datos para fijar el
mapa. La correspondencia demostrada aporta el candidato estructural; todavía no
constituye por sí sola una explicación observacional del sector oscuro.

\paragraph{Dominio de la bisagra areal–volumétrica y condiciones necesarias de validez}
La banda, la razón areal--volumétrica y la lectura bidireccional pertenecen a
la arquitectura previa. La matriz de incidencia, el descenso del calendario,
el envolvente mínimo, la ley de lazos, la medida de Parry, la acción
cuadrática y el retorno marcado proporcionan su formalización explícita. El
lector intervalar y los controles de profundidad \(108\) certifican el
cálculo. Las incidencias \(N_3,N_9,N_{27},N_{108}\), las identidades
matriciales y los cilindros son identidades internas exactas; el teorema de
Möbius y la razón
\(\pi/\varphi^2\) están
\textsc{demostrados con estructura de partida explícita}; las lecturas
partícula--antipartícula e inercial--gravitatoria constituyen interfaces
físicas posteriores; la extensión al sector oscuro permanece como
hipótesis física no demostrada. El par relacionado por la involución
\(\mathscr R(\pi,\varphi)\leftrightarrow
\mathscr R(\varphi,\pi)\) formaliza exactamente una arquitectura previa. Los
coeficientes \((-22,+15)\) quedan demostrados, por el
\cref{p12-83-c15-s2:thm:selector-persistencia-electron-rev7}, con su estructura de partida
explícita.

Una holonomía positiva da un cilindro. Una transición areal--areal, la
ausencia de una transición cruzada o una cuenta TPK distinta de
\(N_9=3F_{\rm av}\) y \(N_{108}=36F_{\rm av}\) falsan el descenso. La
identificación de \(\mu_{\rm Parry}\) con \(\tau_{468}\) o con
\(\nu_{\rm cyc}\) queda refutada por sus valores distintos. Un espectro de
\(\operatorname{Sym}^2(F_{\rm av})\) distinto de
\(\{\varphi^2,-1,\varphi^{-2}\}\), una cuenta \(108\) distinta de
\((F_{107},F_{109})\), un intervalo cociente no anidado o la introducción
iterativa de \(\pi\) dentro de \(F_{\rm av}\) falsan las formalizaciones
nuevas. Una identificación física universal del cociente sin un mapa de
hojas a observables queda fuera del teorema. El delta de las cuatro reglas
queda cerrado; no reabre \(\pi\), \(\varphi\), \(C^*\) ni el empalme
extensión--ruta--registro--carácter.
